\chapter{導入}
\label{ch:introduction}

本資料は，Wasserstein 距離の定義から出発し，
\textbf{目標（距離性；Villani (2009) Chapter 6）}：
Polish 空間 $\mathcal{X}$ 上で
\[
  1\leq p<\infty
  \quad\Longrightarrow\quad
  \Wass_p\text{ は }\Pp{p}(\mathcal{X})\text{ 上の距離である}
\]
を，最適 coupling の存在定理（Villani (2009) Theorem 4.1）を経て証明する．
その後は $p=2$ に限定し，対角 Gaussian 分布の間の $W_2$ を
平均と標準偏差の Euclid 距離として厳密に計算する．
この閉形式を Gaussian decoder に適用して潜在空間上の
引き戻し Riemann 計量を導き，Christoffel 記号，Riemann 曲率，
断面曲率へ進む．最終目標は，decoder が潜在空間に作る曲率を
実際に計算できる形にすることである．
主たる文献は Villani
\emph{Optimal Transport: Old and New}
（Grundlehren der mathematischen Wissenschaften 338, Springer, 2009）であり，
定義・記法・証明はこれに従う．
日本語の解説としては
桑江ほか『最適輸送理論とリッチ曲率』
（Encounter with Mathematics 第63回，中央大学，2015）の §1.4 を参照し，
測度分解定理・接着補題などの用語はこれに合わせた．

\section*{本資料の構成}

\begin{itemize}
  \item 第\ref{ch:wasserstein-metrics}章：
    $\Pp{p}(\mathcal{X})$，coupling，$\Wass_p$ を定義し，
    最適 coupling の存在定理を経て目標を証明する．
  \item 第\ref{ch:gaussian-w2}章：
    $p=2$，$\mathcal{X}=\R^n$ に限定し，1次元および対角 Gaussian
    分布の間の $W_2$ の閉形式を最適 coupling とともに証明する．
    Gaussian decoder が潜在空間に定める引き戻し距離を得る．
  \item 第\ref{ch:latent-curvature}章：
    decoder の平均・標準偏差をまとめた写像の微分から
    引き戻し計量テンソルを導く．その正定値性を確認した後，
    Christoffel 記号と曲率を定義し，2次元の具体例で
    Gaussian 曲率を最後まで計算する．
  \item 付録A：
    距離空間と測度論の準備．
    本編で用いる定義と補題をすべて揃え，
    証明なしで認める標準定理（単調収束定理・一致定理・分解定理など）は
    ステートメントを明示して文献を挙げる．
\end{itemize}

\section*{主な記号}

\begin{itemize}
  \item $\Prob(\mathcal{X})$：$\mathcal{X}$ 上の Borel 確率測度全体．
  \item $\Pp{p}(\mathcal{X})$：有限 $p$ 次モーメントをもつ確率測度全体
    （定義~\ref{def:wp-p-space}）．
  \item $\Couplings(\mu,\nu)$：$\mu,\nu$ の coupling 全体
    （定義~\ref{def:wp-coupling}）．
  \item $\Wass_p(\mu,\nu)$：Wasserstein 距離（定義~\ref{def:wp-wasserstein}）．
  \item $\mathbf{1}_A$：指示関数（定義~\ref{def:meas-indicator}）．
  \item $\norm{f}_{L^p(\mu)}$：$L^p$ ノルム（定義~\ref{def:meas-lp-norm}）．
  \item $\mu\circ T^{-1}$：写像 $T$ による $\mu$ の像測度（定義~\ref{def:meas-pushforward}）．
  \item $\mu\otimes\nu$：直積測度（定理~\ref{thm:meas-product}）．
  \item $\mathcal{N}(m,\operatorname{diag}(\sigma^2))$：
    平均 $m$，対角共分散 $\operatorname{diag}(\sigma^2)$ の
    Gaussian 分布（定義~\ref{def:gw-diagonal-gaussian}）．
  \item $G(z)$：Gaussian decoder が潜在空間に定める
    引き戻し計量テンソル（定義~\ref{def:lc-pullback-metric}）．
\end{itemize}
